\section{Introduction}
\label{sec:intro}

Partial differential equations (PDEs)
are fundamental to modeling physical phenomena across fluid dynamics,
solid mechanics, and thermodynamics.
Solving these PDEs is a critical step in understanding physical mechanisms and
guiding engineering design.
Traditional numerical methods such as finite difference, finite element,
and finite volume methods are widely used but require mesh generation and
iterative solvers that incur significant computational cost,
especially when boundary conditions, material parameters,
or geometric configurations change.

In recent years, physics-informed neural networks (PINNs)
have emerged as a mesh-free alternative for solving PDEs by embedding the
differential operator and boundary conditions directly into the
training objective~\cite{raissi2019pinn}.
However, standard PINNs are typically trained for a single set of physical
parameters; any change to the parameters requires retraining from scratch.
Transfer learning offers a way to mitigate this cost by initializing the target
model with weights from a pretrained source model~\cite{wang2023expert}.

Parameter-efficient fine-tuning (PEFT)
methods such as LoRA~\cite{hu2022lora} and Adapter~\cite{houlsby2019adapter}
have
shown
promise in NLP for adapting large models with minimal parameter changes.
Whether these techniques can effectively adapt PINNs across PDE parameter
shifts, where the loss landscape involves higher-order derivatives
through automatic differentiation,
remains an open question~\cite{biderman2024lora}.
